\section{Related Work}
\label{sec:related}
The performance of \SMT{} solvers is a critical bottleneck in areas like symbolic execution~\cite{csmith, klee, sen2005cute}, especially when faced with complex, non-linear constraints~\cite{DBLP:conf/tacas/BischoffL21}. Prior work has tackled this challenge from multiple angles. We organize the literature into two primary layers---solver-level enhancement and constraint-level simplification---and position \tool as a policy-learning paradigm for \emph{pre-solver} constraint simplification via variable-level semantic concretization.


\subsection{Solver-Level Enhancement}
Solver enhancement approaches modify SMT solvers directly. Widely used SMT backends such as STP~\cite{stp2007}, Z3~\cite{demoura2008z3}, CVC5~\cite{cvc5}, MathSAT5~\cite{mathsat5}, Boolector~\cite{boolector2018}, and BVParti~\cite{bvparti2025} provide a foundation for solver development. Algorithm selection methods~\cite{DBLP:conf/tacas/BischoffL21, scott2023algorithm} use ML to optimize solver configurations and predict solving times; online portfolio selection such as MedleySolver~\cite{pimpalkhare2021medleysolver} further schedules backends adaptively. Internal modifications include neural branching heuristics~\cite{kurin2020improving, wang2021neurocomb} and distributed architectures~\cite{zhao2024distributed} that partition constraints across multiple instances. This solver-level line of work reduces the effective search space via partitioning; by contrast, \tool reduces search space pre-solver through variable-level semantic concretization without touching solver internals. Beyond this, theory arbitrage transforms unbounded constraints into bounded theories to accelerate solving~\cite{mikek2024sta}. Neural solvers~\cite{DBLP:conf/iclr/SelsamB19, DBLP:conf/nips/AbreuGMM21} attempt complete replacement but struggle with complex problems. These approaches are orthogonal to \tool---they build better solvers, whereas \tool reduces difficulty pre-solver without modifying backends and can be composed with them.

In parallel, LLM--solver integration has emerged within this axis, such as Logic-LM~\cite{pan2023logiclm} for integrating LLMs with logical solvers and MCP-Solver~\cite{szeider2025mcp} for constraint programming. These methods strengthen backends or interfaces but, unlike \tool, do not perform variable-level, semantics-aware concretization prior to solving.

\subsection{Constraint-Level Simplification}
Constraint simplification approaches use AI to transform constraints before solving. Deep learning methods~\cite{wenConstraintSolvingDeep2020, yangExploringStrategiesGuiding2024} apply neural networks for constraint prediction and ML-guided symbolic analysis. Strategy synthesis approaches~\cite{chen2021synthesize} automatically generate solving strategies for symbolic execution, while time prediction methods~\cite{luoBoostingSymbolicExecutionTime2021, thobenTimeoutPredictionSoftware2023} use ML to predict constraint solving complexity. Rewrite-based simplification via equality saturation (\eg, egg~\cite{willsey2021egg}) emphasizes canonicalization rather than goal-directed, variable-level concretization. Specialized approaches target specific constraint types~\cite{shuaiTypeIntervalAware2021, liuOptimalRefinementbasedArray2022}, while hybrid methods~\cite{muduliSatisfiabilityModuloFuzzing2022} combine SMT solving with fuzzing techniques. Recent work learns patterns to bias symbolic execution toward vulnerable states~\cite{ruaro2021syml}, and further converts SMT formulas to effective test cases to boost testing pipelines~\cite{zhang2024smt2test}. The latest developments include concrete constraint guidance for symbolic execution~\cite{concrete2024icse} and specialized vulnerability signature generation using static analysis guidance~\cite{zhang2024stase}; \tool instead reduces solver difficulty via variable-level, semantics-aware concretization and composes with these efforts. Unlike equality saturation or handcrafted rewrite rules, \tool performs goal-directed, per-variable concretization guided by learned, semantics-aware policies.

At the constraint level, LLM-based methods have been considered for adjacent tasks such as translating natural language to logical formulas (SatLM~\cite{ye2023satlm}); however, they generally do not learn variable-level, semantics-aware concretization policies for SMT formulas, which is the focus of \tool.
 Taken together, existing work optimizes solvers or applies fixed rewrite strategies; \tool instead offers a cross-cutting, pre-solver strategy that learns semantics-aware, variable-level concretization to reduce solver difficulty.%
\tool operates on a fundamentally different principle from existing constraint simplification approaches. Unlike Wen et al.~\cite{wenConstraintSolvingDeep2020} who use deep learning for general constraint solving, or Yang et al.~\cite{yangExploringStrategiesGuiding2024} who apply ML prediction for symbolic analysis guidance, \tool combines \RL{} and \LLM{} for strategic variable concretization and learns a fine-grained, variable-level policy that actively reduces problem complexity through semantics-aware reasoning. This synergy---where the \RL{} agent provides long-term strategy and the \LLM{} offers semantic intuition for each step---distinguishes \tool from approaches that focus on high-level tactic selection or time prediction~\cite{chakraborty2023ranking, DBLP:conf/icml/LuoCHSZL23, luoBoostingSymbolicExecutionTime2021} and from LLM-based interfaces that do not provide learned, variable-level, semantics-aware concretization for SMT formulas. Overall, LLMs have been applied to formal reasoning and solver integration~\cite{kamath2023finding, chakraborty2023ranking, ye2023satlm, pan2023logiclm, szeider2025mcp}. These techniques largely follow guess-and-verify paradigms with scalability limitations and are complementary to \tool.
